\chapter{Extensiones y límites del algoritmo}
\label{cap:extensiones}

La forma canónica \eqref{eq:canon-v}--\eqref{eq:canon-Y} parece restrictiva.
En la práctica admite mucho más de lo que aparenta, con ajustes menores. Este
capítulo recorre los casos que uno se encuentra al modelar y dice, para cada
uno, si el algoritmo funciona tal cual, si hay que generalizarlo o si hay que
buscar otro método.

\section{Casos que entran sin modificar nada}

\paragraph{Adelantos del insumo} Si en lugar de $X_t$ aparece $X_{t+1}$ en el
paso hacia atrás, el algoritmo funciona sin cambios. El lema
\ref{lema:politicas} sigue valiendo, porque la iteración hacia atrás ya
incorpora el efecto de los insumos futuros a través de la continuación.

\paragraph{Momentos de orden superior} Si uno define el resultado individual como
$y^k$ en lugar de $y$, el agregado es el $k$-ésimo momento no centrado. Con dos
bloques (uno para $y$, otro para $y^2$) y un bloque simple que los combine, se
obtiene el jacobiano de la varianza. Lo mismo vale para un índice de precios CES
o para el coeficiente de variación del consumo.

\paragraph{Representaciones no basadas en grillas} Nada en el algoritmo exige
que $v$ sea un vector de valores sobre una grilla. Puede ser un vector de
coeficientes de Chebyshev o de un spline. Lo que sí tiene que ser una
distribución sobre puntos es $D$ --- y aun eso se relaja en la sección
siguiente.

\paragraph{Elección discreta con choques de gusto} Un modelo con elección
discreta (ajustar o no un bien durable, cambiar o no el precio) tiene políticas
discontinuas, y la perturbación de primer orden sobre una grilla da cero: ningún
punto de la grilla cruza el umbral ante un choque infinitesimal. La solución
estándar es añadir choques de gusto iid --- típicamente de valor extremo, que dan
probabilidades logit --- de modo que la probabilidad de cada opción varía de
forma continua con el estado. Hecho eso, el bloque entra en la forma canónica
tomando $D_t$ como la distribución \emph{antes} de la realización del choque de
gusto, y $v$, $\Lambda$, $y$ como valores esperados sobre ese choque.

\paragraph{Distribución endógena dentro del período} ¿Qué pasa si un choque en
$t$ afecta la posición de los agentes en $t$ --- por ejemplo, una ganancia de
capital sobre la riqueza, o una probabilidad de desempleo que se realiza en $t$?
La regla es: mantener $D_t$ como la distribución \emph{previa} a cualquier evento
de la fecha $t$, e incorporar esos eventos dentro de $v$, $\Lambda$ e $y$.

\begin{trampa}
\textbf{Rezagos del insumo.} Si en el paso hacia atrás aparece $X_{t-1}$, el
lema \ref{lema:politicas} \textbf{falla}: ahora el comportamiento en $t$ depende
del pasado y no solo de la distancia al choque futuro. La solución limpia no es
parchear el algoritmo, sino sacar el rezago del bloque heterogéneo: definir una
variable nueva $\tilde X_t \equiv X_{t-1}$ como producto de un bloque simple
--- que sí admite rezagos sin problema --- y pasarle $\tilde X_t$
contemporáneo al bloque heterogéneo. En el modelo de Krusell--Smith del
capítulo \ref{cap:jacobiano}, el rezago del capital vive en el bloque de
empresas por exactamente esta razón.
\end{trampa}

\section{La generalización: productos y transiciones no lineales}

Algunos modelos necesitan reemplazar \eqref{eq:canon-D} y \eqref{eq:canon-Y} por
funciones generales de la distribución:
\begin{align}
D_{t+1} &= \mathcal{D}(v_{t+1}, X_t, D_t),\\
Y_t &= \mathcal{Y}(v_{t+1}, X_t, D_t).
\end{align}

\begin{proposicion}[\ABRS, proposición 2]
\label{prop:general}
Con esas dos ecuaciones generales, la proposición \ref{prop:principal} sigue
valiendo si se redefine el vector de expectativa como
\begin{equation}
\E_t \equiv (\mathcal{D}_D')^{\,t}\,\mathcal{Y}_D',
\label{eq:E-general}
\end{equation}
donde $\mathcal{D}_D$ es el jacobiano $n_g\times n_g$ de $\mathcal{D}$ respecto de
$D$ y $\mathcal{Y}_D$ el gradiente de $\mathcal{Y}$ respecto de $D$, ambos
evaluados en el estado estacionario.
\end{proposicion}

La demostración es la misma del lema \ref{lema:fake} reemplazando $\Lambda_\sst'$
por $\mathcal{D}_D$ y $y_\sst'$ por $\mathcal{Y}_D$. El contenido económico es
transparente: $\mathcal{D}_D$ es el ``operador de propagación'' de una
perturbación distributiva y $\mathcal{Y}_D$ es el ``operador de lectura'' del
agregado. En el caso lineal original, esos dos operadores son $\Lambda_\sst'$ e
$y_\sst'$.

Esto abre tres casos útiles.

\paragraph{Entrada y salida de agentes} Si los agentes salen con probabilidad
$1-\varsigma$ y entran con una distribución fija $\Psi$,
\begin{equation}
D_{t+1} = \varsigma\, \Lambda(v_{t+1},X_t)'\,D_t + (1-\varsigma)\,\Psi .
\end{equation}
Entonces $\mathcal{D}_D = \varsigma\,\Lambda'_\sst$ y $\mathcal{Y}_D = y_\sst'$. La
única consecuencia práctica es que los vectores de expectativa se descuentan por
la supervivencia: $\E_t = \varsigma^t \Lambda_\sst^t y_\sst$. Todo lo demás del
algoritmo queda igual. Este es el caso de los modelos de dinámica de empresas y
de los modelos bancarios tipo Gertler--Karadi, y es el que usamos en la parte V.

\paragraph{Distribuciones paramétricas} Si $D_t$ es un vector de parámetros que
describe una densidad en lugar de masas sobre puntos, basta con poder calcular
$\mathcal{D}_D$ y $\mathcal{Y}_D$, típicamente por diferenciación numérica.

\paragraph{Momentos no polinómicos} Un cuantil, un índice de Gini, una razón
entre percentiles: todos son funciones no lineales de $D$, y el algoritmo los
admite vía $\mathcal{Y}_D$. Con una distribución discretizada el cuantil es una
función escalonada y su derivada no sirve; la solución es interpolarla
linealmente entre puntos de masa, lo que da un gradiente analítico sencillo.

\section{Lo que no entra}

Hay una restricción que la generalización no levanta: la proposición
\ref{prop:general} deja intacta la ecuación \eqref{eq:canon-v}. La variable de
continuación $v_t$ \textbf{no puede depender de $D_t$}.

Económicamente, eso excluye todo modelo en el que a un agente le importe
directamente la distribución futura de otros agentes, de una forma que no pase
por un precio o agregado $X_t$. Los ejemplos de la literatura son:

\begin{itemize}[leftmargin=1.4em]
\item modelos de generaciones superpuestas con herencias endógenas recibidas a
mitad de la vida;
\item modelos de búsqueda laboral con publicación de salarios o negociación
individual, donde importa la distribución de ofertas;
\item modelos de búsqueda monetaria donde importa la distribución de saldos
de efectivo de las contrapartes.
\end{itemize}

\begin{tutesis}
Vale la pena chequear esta condición en tu modelo, porque está más cerca de lo
que parece. Si el mercado interbancario se modelara con \emph{búsqueda y
negociación bilateral} al estilo de Afonso y Lagos (2015), a cada banco le
importaría la distribución de posiciones de liquidez de sus contrapartes
potenciales, y la condición fallaría. Tu plan de trabajo evita ese problema
\emph{por construcción}: la regla proporcional $\psi_k = \tau_k \min(1, P/Q)$
hace que toda la información del mercado llegue al banco a través de dos
escalares agregados, $P$ y $Q$. Esa simplificación, que el plan justifica como
``una aproximación reducida'', resulta ser además exactamente la condición que
hace aplicable el método. Conviene decirlo explícitamente cuando defiendas esa
elección de modelado: no es solo tratabilidad, es lo que mantiene el modelo
dentro de la clase resoluble.
\end{tutesis}

\section{Varias etapas dentro de un período}

Un caso frecuente en macrofinanzas: dentro del período $t$ pasan varias cosas en
orden --- primero se eligen carteras, después se realiza un choque, después hay
un mercado que redistribuye liquidez, después se paga. Modelar eso metiendo
todo dentro de $v$, $\Lambda$, $y$ se vuelve rápidamente ilegible.

La generalización natural es permitir varias \emph{etapas} por período, cada una
con su propio operador de transición, y encadenarlas. El algoritmo se aplica
tratando la composición de etapas como un único $\mathcal{D}$. En la práctica, si
una de las etapas es un mercado cuyo precio se determina por vaciado, la forma
más simple de proceder es la del capítulo \ref{cap:dag}: declarar ese precio
como incógnita agregada y su condición de vaciado como objetivo, y dejar que el
sistema de equilibrio general lo resuelva simultáneamente.
